\originalchapter{开放系统与光学Bloch方程}\label{ch:open}
\originalsection{环境迹与Markov近似}
原子与辐射场的联合状态仍按幺正定律演化, 但只观察原子时必须取部分迹. 从初始乘积态$\rho_S\otimes\rho_E$和联合幺正$U$得$\rho_S'=\sum_\mu K_\mu\rho_SK_\mu^\dagger$, $K_\mu=\bra\mu U\ket{0_E}$, 且$\sum K_\mu^\dagger K_\mu=I$. 这说明耗散并非把Schrödinger方程任意改成非幺正, 而是忽略了可携带信息的环境.

对弱耦合宽带真空, 二阶相互作用图积分含环境相关函数$C(\tau)$. 当其衰减时间$\tau_E$远小于原子的布居与相干变化时间, 可把$\rho_S(t-\tau)$换为$\rho_S(t)$并把积分上限延到无穷. 再分离不同Bohr频率的快振荡项, 得到保持完全正性的Lindblad结构. 对二能级自发发射,
\begin{equation}\dot\rho=-\frac\ii\hbar[\widetilde H,\rho]+\Gamma\left(\sigma_-\rho\sigma_+-\frac12\{\sigma_+\sigma_-,\rho\}\right).\label{eq:master}\end{equation}
环境记忆长、强耦合或接近不同跃迁简并时, 上述近似需要重新检查. 不能仅因公式写得简洁就认定所有实验都Markov.

\originalsection{稳态、饱和与线宽}
在$\ket e,\ket g$基中, 令$u=2\Re\rho_{eg},v=-2\Im\rho_{eg},w=\rho_{ee}-\rho_{gg}$, 由式\eqref{eq:master}逐项求矩阵元得
\[\dot u=\Delta v-\frac\Gamma2u,\quad\dot v=-\Delta u-\Omega w-\frac\Gamma2v,\quad\dot w=\Omega v-\Gamma(w+1).\]
置三导数为零, 第一式给$u=2\Delta v/\Gamma$, 第三式给$w=\Omega v/\Gamma-1$; 代入第二式可解得
\begin{equation}\rho_{ee}^{\rm ss}=\frac{\Omega^2/4}{\Delta^2+\Gamma^2/4+\Omega^2/2}=\frac{s/2}{1+s+(2\Delta/\Gamma)^2},\qquad s=\frac{2\Omega^2}{\Gamma^2}.\label{eq:ss}\end{equation}
强驱动极限$\rho_{ee}\to1/2$, 因为同一光场既吸收又受激发射. 弱驱动散射率$R=\Gamma\rho_{ee}$, 共振最多$\Gamma/2$. 谱线的角频率全宽为$\Gamma\sqrt{1+s}$, 这就是功率展宽. 闭合偶极跃迁且按本卷$\Omega$定义时, $s=I/I_s$, $I_s=\pi hc\Gamma/(3\lambda^3)$; 角动量分支与非闭合跃迁会改变有效饱和强度.

\originalsection{弛豫、纯退相干与量子轨迹}
无驱动时$\rho_{ee}(t)=\rho_{ee}(0)\e^{-\Gamma t}$,$\rho_{eg}(t)=\rho_{eg}(0)\e^{-\Gamma t/2}$, 故$T_1=1/\Gamma,T_2=2/\Gamma$. 额外相位噪声率$\gamma_\phi$给$1/T_2=1/(2T_1)+\gamma_\phi$. 随机相位$\theta$若均值零、方差$2\gamma_\phi t$, 用高斯特征函数$\langle\e^{\ii\theta}\rangle=\e^{-\gamma_\phi t}$即可得到相干衰减, 而布居不变.

一次短时间监测的跳跃算符为$K_1=\sqrt{\Gamma\dd t}\sigma_-$, 无计数算符为$K_0=I-\ii H\dd t/\hbar-(\Gamma\dd t/2)\ket e\bra e$. 有计数时归一化状态跳到$g$; 无计数时激发幅度相对减小, 因“没有看到光子”也是信息. 把两种结果按概率平均便恢复式\eqref{eq:master}. 不同监测方式可产生不同条件轨迹, 同一未条件化主方程不能唯一推断环境里实际发生哪种测量.
